\part{Results and Evaluation}

\chapter{Results \& Evaluation}
\label{chap:results}

\section{Executive Summary - Three-Pipeline Comparison}

\begin{table}[H]
\centering
\caption{Comparison of Three Pipelines}
\label{tab:pipelines}
\begin{tabular}{|l|l|l|l|l|}
\hline
\textbf{Pipeline} & \textbf{Accuracy} & \textbf{Innovation} & \textbf{Use Case} \\
\hline
Signal-Based (v4) & 100\% (Test) & 4-architecture comparison with v4 optimizations & Clinical ECG machines \\
\hline
Image-to-Signal (U-Net) & MSE: 0.016 & Digitize scanned ECG images for downstream signal models & Old medical records \\
\hline
Image-to-Classification (ResNet) & 82\% (Test) & Direct image classification bypassing digitization & Mobile photo diagnosis \\
\hline
\end{tabular}
\end{table}

\textbf{Literature Benchmarks:}
\begin{itemize}
    \item Signal-based (Kaggle pre-segmented): 100\%
    \item Image-to-Signal (U-Net): MSE 0.016
    \item Image-to-Classification: 82\%
\end{itemize}

\section{Signal-Based Classification (Kaggle v4)}
\textbf{Training Configuration:} 4 models × 30 epochs, AMP enabled, Focal Loss, ReduceLROnPlateau, augmentation.

\subsection{Training Results}
\textbf{Source:} Kaggle kernel \texttt{ecg-research-v4-optimized-training}

\begin{table}[H]
\centering
\caption{Signal-Based Model Performance}
\label{tab:signal_results}
\begin{tabular}{|l|c|c|c|c|c|}
\hline
\textbf{Model} & \textbf{Train Acc} & \textbf{Val Acc} & \textbf{Train F1} & \textbf{Val F1} & \textbf{Epochs to Converge} \\
\hline
ECG\_ResNet & ~100\% & ~100\% & ~1.0 & ~1.0 & ~22/30 \\
\hline
ECG\_LSTM & ~100\% & ~100\% & ~1.0 & ~1.0 & ~22/30 \\
\hline
ECG\_MultiModule & ~100\% & ~100\% & ~1.0 & ~1.0 & ~22/30 \\
\hline
ECG\_PhysioWave & ~100\% & ~100\% & ~1.0 & ~1.0 & ~22/30 \\
\hline
\end{tabular}
\end{table}

\textbf{Observations:}
\begin{itemize}
    \item Pre-segmented Kaggle dataset (each sample = one heartbeat) enables near-perfect classification
    \item All models converge by epoch 22/30
    \item Loss values ($10^{-8}$ to $10^{-10}$) indicate complete convergence on this simplified task
    \item \textbf{Important Note:} This is NOT a medical breakthrough! the dataset is pre-segmented, making classification easier than real-world continuous ECG
\end{itemize}

\subsection{Test Set Results}
\textbf{The test phase was conducted across all four architectures, totaling 120 training epochs to ensure convergence.}

\textbf{Test Performance:}
\begin{itemize}
    \item \textbf{Test Accuracy:} 100\% (pre-segmented MIT-BIH Kaggle CSV)
    \item \textbf{F1-Score (weighted):} 1.00
    \item \textbf{Class-wise Performance:} Focal Loss should improve recall on minority classes (F, S beats)
\end{itemize}

\textbf{Comparison with Literature:}
\begin{itemize}
    \item Nature Medicine (2019): ~90\% recall on \textit{continuous} ECG recordings
    \item Kaggle benchmark: 97-100\% on pre-segmented data (our results align)
    \item \textbf{Conclusion:} Our pipeline works correctly, but real challenge lies in raw ECG recordings
\end{itemize}

\section{Image-to-Signal Pipeline (U-Net)}
\textbf{Source:} Kaggle kernel \texttt{ecg-image-to-signal-v1}.

\textbf{Results:}

Train Loss: 0.016909, Val Loss: 0.016022

\textbf{Use Case:} Digitize scanned ECG images for downstream classification with signal models.

\section{Image-to-Classification (ResNet-50)}
\textbf{Source:} Kaggle kernel \texttt{ecg-image-to-class-v1}.

\textbf{Results:}

\begin{table}[H]
\centering
\caption{Image-to-Classification Results}
\label{tab:image_class_results}
\begin{tabular}{|l|l|l|l|}
\hline
\textbf{Model} & \textbf{F1} & \textbf{Comparison to Signal-Based} \\
\hline
ResNet-50 & 0.82 & Harder than signal-based (100\%) due to spatial encoding loss \\
\hline
\end{tabular}
\end{table}

\textbf{Why lower than signal-based?} Converting 1D→2D→1D loses some precision; image classification is inherently harder.

\section{Dimensionality Reduction Analysis}

\begin{itemize}
    \item \textbf{PCA:} Project 187-dim signals to 2D, visualize 5-class separability
    \item \textbf{t-SNE:} Non-linear dimensionality reduction for cluster visualization
    \item \textbf{Feature Comparison:} Compare learned representations across 4 architectures
\end{itemize}
